% ID: SEL-D0008-BC-E05
% ESTADO: cotejada. Segunda pasada visual completa y borrador compilado revisados.
% Fuente: Selectividad/matematicas_aplicadas/ORD_SUPLENTE_Modelo-4-Enunciados (Mat-Ap-2021).pdf
% SHA-256: 74b067b81be20aa0262a0cc58795f3b25488aa9c749d825b7b33f656dd1ae089
% Página: 2; bloque C; ejercicio original 5.

Una determinada ciudad tiene en la plantilla del ayuntamiento 1000 agentes
de la policía local, 600 bomberos y 400 funcionarios de protección civil.
En esta plantilla, el 42\% de policías, el 20\% de bomberos y el 50\% de
funcionarios de protección civil son mujeres. Se elige una persona al azar
de la plantilla.
\begin{enumerate}[label=\alph*)]
  \item \textbf{(1.5 puntos)} \textquestiondown{}Cuál es la probabilidad de
  que sea mujer?
  \item \textbf{(1 punto)} Si la persona elegida es hombre,
  \textquestiondown{}cuál es la probabilidad de que sea bombero?
\end{enumerate}
